\scanpage{78}
% archive sheet number: 77

\unsure{Questions|Question} sur les topos. (Pour Mlle \unsure{Zunigery|Zunigey|Lucigey} ?).

1) Définition et étude des ``immersions'' de topos. \unsure{ou|de|des} Caractérisation ? Par condition \(f_*\) pl fid i.e. \(f^*f_*\xrightarrow{\sim}\operatorname{id}\) (i.e. \(f^*\) passage au quotient? Dans quelle mesure la \deleted{\unsure{NB|N.B.}} connaissance de \(f:\unsure{\mathcal U|U}\longrightarrow\unsure{\mathcal V|V}\) détermine-t-elle la ``partie loc fermée'' lui donnant naissance ? Plus précisément, si

% slanted
\margin{Théorie des \(f_!\) et \(f^!\)}

% slanted
\margin{\unsure{faux}}

\margin{OK. pour imm. ouverte et imm. fermée.}

\diagram{\begin{tikzcd}[ampersand replacement=\&,row sep=tiny]
\mathcal U\arrow[r,<->,"f"]\arrow[dr,<->,"{f'}"']\&\mathcal V\\
\deleted{\unsure{\mathcal U'|U'}}\arrow[r,draw=none]\&\mathcal V'
\end{tikzcd}}

sont des hom d'imm. de topos de \(=\) \deleted{\unsure{source}} \inserted{but} qui sont ``équivalents à \deleted{\unsure{source}} \inserted{but} fixé'' [i.e. \(\exists\) équivalence

\diagram{\begin{tikzcd}[ampersand replacement=\&,row sep=small]
\mathcal V\arrow[rr,"g"]\arrow[dr,"{f^*}"']\&\&\mathcal V'\arrow[dl,"{f'^*}"]\\
\&\mathcal U\&
\end{tikzcd}}

rendant comm. à isom près le diagramme], alors \(g\) est-elle déterminée à isom unique près ? [Notion de ``monomorphisme essentiel de Topos'' ??].

Caractérisation des imm ouvertes [par \(f^*\) commutant aux \(\underleftarrow{\lim}\) quelc., en plus de \(f_*\) pl fid ?].

\margin{\uline{\unsure{non|voir}}}

Caractérisation des imm fermées [par la condition \(F\longrightarrow f_*f^*F\) un épim pour \unsure{tt|\(H\)} faisceau \uline{abélien} \(F\) -- attention, \unsure{contre|autre} en dehors de ce cas, prendre \(F=\varnothing\) \(\ldots\) -- en plus de \(f_*\) pl fid.] ?

% slanted
\margin{Application \unsure{aux} ``bonnes corresp.'' pour \unsure{commutation} \deleted{\unsure{leurs}} \(H^i\) aux \(\underrightarrow{\lim}\)}

\textcircled{2} Etude des \(f\) tels que \(f^*\) commute aux \(\underleftarrow{\lim}\) quelc. Alors \(f^*\) \deleted{est-il de la forme \(\unsure{g_*}\)} admet-il un adjoint à \uline{gauche} \deleted{\unsure{\(g\)}} \(f_!\) ? Dans quelle mesure le cas du foncteur \unsure{canonique} \(f^*:\deleted{\unsure{V}}\mathcal U\longrightarrow\mathcal U_{/S}\) est-il typique ??

% slanted
\margin{\uline{oui}}

\margin{Il ne l'est pas, cf pts essentiels. Regarder le cas des commutation aux \uline{Hom} relatifs. \(\ldots\)}
